%% Chapter 8: Discussion and Conclusion

\chapter{Discussion and Conclusion}
\label{chap:discussion_conclusion}
This chapter wraps up the thesis. We begin with a synthesis of our findings in light of the motivating research question. Then we look at the contributions of the thesis relative to existing literature, and discuss the gap between the ABMs of Chapter~\ref{chap:abm} and the analytical modelling of subsequent chapters. We end by discussing the implications for SMU--WB research, along with limitations and directions for future work.

%% Section 8.1: Synthesis of Findings
\section{Synthesis of Findings}
\label{sec:discussion_conclusion:synthesis}
We now return to the motivating research question of Section~\ref{sec:intro:motivation}:
\begin{quote}
\emph{What is the causal influence of social media use on well-being?}
\end{quote}

The answer is that the causal influence of SMU on WB remains unclear. This thesis has not tried to answer the question directly, but instead shows that omission of network effects might be a key reason why the existing evidence cannot answer it. We developed network-aware models for future empirical estimation.

In Chapter~\ref{chap:introduction} we argued that the field of research tackling this question is affected by the methodological problem of not modelling network effects. This is because the standard statistical methods of the field, such as cross-sectional OLS, panel data with fixed effects, and randomised controlled trials (RCTs), fail when applied to networked data. In particular, we identified three statistical problems that arise when network effects are neglected: \emph{endogeneity} (peer outcomes enter the DGP but not the regression), \emph{relative effects} (differences between individual and peer behaviour have an effect on WB), and \emph{spillovers} (changes in one unit propagate through the network, breaking independence). The endogeneity problem relates to the \emph{reflection problem} of \textcite[Sec.~2]{Manski1993}, which is the problem of separately identifying endogenous, contextual, and correlated peer effects within a linear-in-means formulation.

We now give a synthesis of how the thesis arc follows these problems, from the problem diagnosis of Chapter~\ref{chap:abm} to the solutions developed in Chapters~\ref{chap:model}--\ref{chap:simultaneous}.

\subsection{Diagnosis: Standard Methods Fail Under Network Effects}

We begin by recounting the diagnosis of Chapter~\ref{chap:abm}, which applied ABMs to networked data. We defined three scenarios: a reference scenario, a scenario with the CoMO effect using a global mean, and a Network \& CoMO scenario based on peer reference. In the reference scenario, standard methods recovered the true effect. However, when the true DGP was networked, simulations showed attenuation of the true effect across the board. In the global CoMO scenario the methods exhibited relative biases of roughly $20\%$--$26\%$, with the cross-sectional OLS and RCT biases increasing further in the Network \& CoMO scenario (see Table~\ref{tab:abm_summary}). The most noticeable result was that RCTs, often regarded as the gold standard of causal inference, were the \emph{most} biased method in our simulations, with a relative bias of $31.9\%$ in the Network \& CoMO scenario (but note that this is subject to the $n$/density caveat from Section~\ref{sec:abm:methods}). The explanation is the independence required by the SUTVA Assumption~\ref{ass:sutva}. When we treat some agents with less SMU, this has a spillover effect on others through lowering their peer reference level. This attenuates the contrast between the groups, biasing estimates downwards.

In addition to showing that standard methods are biased under network effects, we used the complexity allowed for by ABMs to demonstrate Simpson's Paradox in Section~\ref{sec:abm:simpsons}. Implementing a delayed-decline mechanism allowed us to show how a DGP with a negative influence of SMU on WB could still lead to the paradox: cross-sectional regression at any single time point shows a positive association between SMU and WB, while at a population level, average SMU is increasing whereas average WB is declining. This shows the difficulty research in this field faces, as someone applying the seemingly reasonable procedure of cross-sectional regression might actually retrieve estimates with the opposite sign to the true effect. As many studies in the SMU--WB field rely on exactly such single-survey regressions, these results, together with the bias results, suggest a potential explanation for the conflicting estimates in the literature: substantial harms in \textcite{Twenge2018} and \textcite{Haidt2024}, and negligible effects in \textcite{Orben2019}. This divergence in results might be due to methodological limitations of methods not taking network effects into account.

Our finding of bias towards zero applies most strongly to studies that report small effects. If the underlying effect is actually large, but this type of bias or attenuation is present, we would expect to find small effects. We note that our simulations provide evidence for the plausible direction of bias in general for studies not taking network effects into account, which does not necessarily apply to any individual study.

\subsection{Solution: A Network-Aware Model with Formal Estimation}

We now move on to the solution: explicit modelling of network effects. In Chapter~\ref{chap:model} we introduced the recursive CoMO model, which lets peers influence the individual through three channels: peer conformity in SMU through $\beta_1$ in the SMU equation, and the CoMO effect and endogenous peer-WB spillover through $\gamma_2$ and $\gamma_3$ in the WB equation. This model addresses the three statistical problems caused by neglecting network effects. We recognise $\mat{G}\vect{y}$ as endogenous, and handle it accordingly when developing estimators. Relative peer effects are formalised and captured through the asymmetric CoMO term. Spillovers are made explicit in the peer-outcome term of the reduced form, and accounted for in estimation. The recursive structure makes the equations separately estimable under Assumption~\ref{ass:error_indep}, which simplifies estimation.

The CoMO term is non-linear, but does not cause problems in estimation. Instead, it provides us with the additional excluded instrument $\mat{G}\vect{c}$, which a standard SAR model does not usually have. We can use this as an instrument because the max operator makes $\vect{c}$ generically linearly independent of $\vect{1}$ and $\vect{s}$. This is formalised as a full-column-rank condition on the design matrix $\mat{X} = [\vect{1}, \vect{s}, \vect{c}]$. In Proposition~\ref{prop:wb_identification}, we use this instrument together with $\mat{G}^2\vect{s}$ to show that the WB structural parameters are identified via the IV route under the \textcite[Prop.~1]{Bramoulle2009} condition that $\{\mat{I}, \mat{G}, \mat{G}^2\}$ are linearly independent. We also show how the Gaussian form of the likelihood can identify the full WB parameter vector (Proposition~\ref{prop:wb_ml_identification}).

We developed both ML and 2SLS estimators for the model in Chapter~\ref{chap:estimation}, along with the analytical Hessians. For the SMU equation, we developed only the ML estimator, as the lack of exogenous covariates makes 2SLS unavailable. The estimators have different properties. The ML estimator requires normality, but is asymptotically efficient when this holds. The 2SLS estimator is more robust to distributional misspecification, but is generally less efficient. Applying both estimators is informative: agreement between them strengthens confidence in estimates, while disagreement requires consideration of whether there are weak instruments or an error in the specification.

\subsection{Validation and Extension}

The developed estimators were validated through Monte Carlo simulations in Chapter~\ref{chap:monte_carlo}. The main result is the success of both estimators in recovering true parameters. They have low bias, standard errors declining at an expected rate of $1/\sqrt{n}$, and roughly nominal confidence interval coverage. The analytical Hessians were verified to match closely with numerical estimates. The simulations also confirmed the theoretical properties of the two estimators, showing the increased efficiency of the ML estimator across the simulation studies. The relative efficiencies ranged from $1.2\times$ to $5\times$ across the parameters, with $\gamma_3$ having the largest gap, likely due to noise caused by weak instruments. When normality seems reasonable, these results clearly favour the ML estimator. Another set of notable results with practical consequences comes from the density study. These showed that as the density of the network increased, instrument strength weakened. From the least to most dense networks tested, the first-stage $F$-statistic fell from roughly 15 to 3.5, and RMSE for $\gamma_3$ went from $0.305$ to $2.759$ for the 2SLS estimator. The ML estimator performance also degraded, but at a modest rate (see Section~\ref{sec:mc:density}). Thus, for networks of modest to high density, which might include networks most likely to meet the data requirements of Section~\ref{sec:discussion_conclusion:data_requirements}, the ML estimator should be the primary estimator, with the 2SLS serving as a check of robustness.

Lastly, we extended the recursive CoMO model to the simultaneous CoMO model in Chapter~\ref{chap:simultaneous} by relaxing the recursive assumption. We allowed for feedback from WB to SMU through an additional parameter $\beta_2$ in the SMU equation. This led to a system without a closed-form reduced form, which required showing the existence of a unique equilibrium. We showed this in Proposition~\ref{prop:contraction}, together with the conditions under which fixed-point iteration converges to said equilibrium. We then developed an FIML estimator for the model, along with a Schur-complement simplification that exploits sparsity for computational efficiency. We provided an initial validation of the estimator with Monte Carlo simulations, which showed low bias and well-calibrated standard errors. Additionally, as the recursive CoMO model is nested in the simultaneous CoMO model as the special case of $\beta_2 = 0$, we applied a Wald test for detecting the need for the simultaneous CoMO model under a simultaneous DGP\@. The test size was validated, with test power being low to moderate in our simulation study.

%% Section 8.2: Contributions in Context
\section{Contributions in Context}
\label{sec:discussion_conclusion:contributions}
In Section~\ref{sec:intro:contributions} we outlined the six novel contributions made in this thesis. This section discusses these in light of the wider literature.

\textbf{(i) The ABM demonstration of bias and Simpson's Paradox} connects to the broader literature on interference and the SUTVA assumption \parencite[Sec.~2.1]{HudgensHalloran2008}, as well as to the tradition of using ABMs to understand social phenomena \parencite{Macy2002}. We are not aware of any prior studies quantifying bias in the SMU--WB setting under the CoMO mechanism. This thesis thus contributes by placing the SMU--WB question both in the frame of the interference literature and in the methodological tradition of understanding through simulated bottom-up dynamics.

\textbf{(ii) The recursive CoMO model} builds on psychological literature that has discussed social comparison since \textcite{Festinger1954}. Reframing the FOMO concept of \textcite[p.~1841]{Przybylski2013}, \textcite[Eqs.~1--2]{DahlFoldnesGronneberg2024} introduced the CoMO concept in the SMU--WB context and used a stylised linear mediation model to illustrate some of its consequences, with \textcite{DahlFoldnesGronneberg2026} extending the framing for policy applications. The contribution made in this thesis is the development of the concept into an identified component of a SAR model. We adapt the \textcite{Bramoulle2009} identification argument to show identification of the WB equation in the presence of the CoMO term.

\textbf{(iii) and (iv) The ML and 2SLS estimators} are developed in line with existing theory of maximum likelihood for SAR models of \textcite{Lee2004}, and the \textcite{Bramoulle2009} network instrument approach for 2SLS estimation. For the recursive CoMO model, our contribution is to adapt this theory to the two-equation structure with the non-linear CoMO term. Checking model robustness through applying both ML and 2SLS estimators is a known practice in spatial econometrics \parencite{LeSage2009}.

\textbf{(v) The Monte Carlo validation} of the estimators is applied to sample sizes relevant for empirical application. In particular, the density sensitivity study of Section~\ref{sec:mc:density} contributes an important insight, which is that 2SLS becomes unreliable when networks reach a certain density, making the ML estimator preferable for relevant applications.

\textbf{(vi) The simultaneous CoMO model} moves further away from the standard linear-in-means systems commonly treated in spatial econometrics, contributing an explicit uniqueness result based on Banach's fixed-point theorem. The previous SMU--WB literature does not provide \emph{a priori} knowledge of whether feedback from WB to SMU is negligible. Thus, this extension contributes by broadening the applicability of the developed framework, without sacrificing formal inference.

The contributions of this thesis are methodological, and application of the methods to real SMU--WB data remains future work (Section~\ref{sec:discussion_conclusion:limitations_future}).

%% Section 8.3: Bridging ABMs and Analytical Methods
\section{Bridging ABMs and Analytical Methods}
\label{sec:discussion_conclusion:gap}
A central approach of this thesis is the complementary application of ABM simulations and analytical inference. This section bridges the two, discussing their strengths and weaknesses.

ABMs can represent complexity in a way analytical methods cannot \parencite[Ch.~1]{Epstein2006}. In Chapter~\ref{chap:abm}, we set up heterogeneous agents with varying parameters and gave them learning rules, mean-reverting behaviour, delayed effects, and network-dependent CoMO\@. The simulations then evolved in complex ways according to these rules over time. This level of complexity in the DGP is intractable to model with analytical approaches. However, the real world might be closer to the complexity of ABMs than to that of analytical methods. Thus, ABMs can give insight into real-world dynamics that would otherwise be out of reach. In our case they are our diagnostic tool, as they show how network effects bias estimates of standard methods, as well as how such effects can lead to counterintuitive emergent results such as Simpson's Paradox.

Analytical methods allow for formal inference, letting us quantify uncertainty in a way not available with ABMs. For our recursive CoMO model of Chapters~\ref{chap:model}--\ref{chap:estimation}, this means that we have the properties of consistency and asymptotic normality, and we have available to us standard errors, confidence intervals, hypothesis tests, as well as conditions on identifiability. These allow us to draw more credible causal conclusions from a single observed dataset when the model assumptions hold. While the ABMs can diagnose bias issues under network effects, the analytical estimator gives us the underlying parameter estimates, together with measures of how precisely they have been estimated.

The gap between the two approaches is found in the complexity analytical approaches cannot handle, but which ABMs can. In our approach, the clearest gap is in dealing with complex temporal effects of agentic behaviour. For example, the CoMO mechanism is applied in the ABMs in a way that allows it to interact with learning and adaptation, generating feedback loops and other temporal dynamics. In our recursive CoMO model, the mechanism is static, capturing the CoMO an individual experiences at a single cross-section. We need this simplification to allow for identification and estimation, but it is arguably a strong assumption. Compared to the standard methods shown to be biased, our recursive CoMO model, and especially the simultaneous CoMO model of Chapter~\ref{chap:simultaneous}, allow for analytical modelling of peer effects. Thus, they can be seen as a start to narrowing the gap between ABMs and analytical approaches. However, the gap remains wide, and requires substantial further work to be closed.

We do not apply our analytical estimators directly to the ABM-generated data of Chapter~\ref{chap:abm}. This is because the DGP used for the ABMs is outside the parametric family of the recursive CoMO model, with temporal structure, learning rules, and heterogeneity in agents that break the model assumptions found in Section~\ref{sec:model:assumptions}. Developing analytical models fit for the DGP used in our ABMs requires substantial future work. In this thesis, we are content with validating our analytical estimators on the DGP they were designed for.

In our thesis, both approaches were used not only to draw the line from diagnosing with ABMs to analytical solutions, but also to illustrate that the gap persists. This strategy of using ABMs for diagnosis and analytical methods for inference could be applied more generally to areas where network effects are present, or where other effects require new estimation methods.

%% Section 8.4: Implications for SMU--WB Research
\section{Implications for SMU--WB Research}
\label{sec:discussion_conclusion:implications}
The finding that network effects can lead to bias, or even the wrong sign, in estimates of the relationship between SMU and WB when using standard methods gives a new reading of existing studies. The social nature of social media makes the presence of such effects plausible. Thus, the conflicting evidence in the current literature should not be read as conclusive evidence of either negligible or strong effects. Instead, it might be better to interpret this as an indication that network effects (or other unmodelled effects) are present, and that this leads to biased and inconsistent estimates.

The simulations in this thesis use parameters and network topologies that are not based on real data. Hence, we should be careful about over-interpreting the results. The bias magnitudes and sign-shift in the Simpson's Paradox illustration should therefore be taken as evidence that such effects are \emph{possible}, but not that these exact magnitudes or effects are present. Other unknown factors might be present, so the takeaway is the need for a careful evaluation to ensure appropriate modelling and interpretation of estimates.

The current research should also take into account the possibility that the CoMO mechanism is empirically relevant. While the mechanism has not yet been validated with real data, a true $\gamma_2 < 0$ in the population would require the field to adapt its approach. In that case, investigating the SMU--WB relationship requires a focus not only on the absolute usage of each individual, but also on how this compares to the use of their peers. For example, someone using social media for two hours per day could have higher well-being when their peers use it for one hour than if their peers use it for four, and we would need information on the peer usage to know the full story. This has consequences for the type of intervention designs future studies should adopt. Changing individual-level SMU might be counterproductive if peer SMU stays the same, as the increased gap with peers would reduce WB through the CoMO mechanism. Such counterproductive effects can be alleviated by introducing interventions on groups of peers or larger communities simultaneously, avoiding the creation of such gaps \parencite[Sec.~2.2]{HudgensHalloran2008}.

The models developed in this thesis would best contribute to the literature through an application to real data combining SMU, WB, and network measurements for the same individuals. Doing this using cluster-randomised experimental designs \parencite[Sec.~2.2]{HudgensHalloran2008}, where whole peer groups are randomised together, would mitigate the SUTVA violations documented in Chapter~\ref{chap:abm}. This would provide the field with a firmer causal basis than the literature currently offers.

%% Section 8.5: Applying the Framework in Practice
\section{Applying the Framework in Practice}
\label{sec:discussion_conclusion:practice}
Using the methodological framework developed in this thesis requires data in a specific form. This section outlines these requirements, and provides practical guidance based on the Monte Carlo evidence of Chapter~\ref{chap:monte_carlo} on model estimation.

\subsection{Data Requirements}
\label{sec:discussion_conclusion:data_requirements}
The list below covers the three types of information our model estimation procedure expects. These should ideally be captured at the same time point to align with the assumptions of our cross-sectional model.

\begin{enumerate}
    \item \textbf{Network data.} For the network data we need to observe the adjacency matrix $\mat{G}$. Thus, we ideally need a complete mapping of the social connections present in our sample. If this is not at hand, we need to have available a reliable approximation. Use of partial network data may bias both ML and 2SLS estimates. We also need this observed $\mat{G}$ to meet a condition for identification. For IV identification and 2SLS estimation, linear independence of $\{\mat{I}, \mat{G}, \mat{G}^2\}$ is required (Proposition~\ref{prop:wb_identification}). If ML identification and estimation are used, we need the linear independence of $\{\mat{I}, \mat{G} + \mat{G}^\top, \mat{G}^\top\mat{G}\}$ (Propositions~\ref{prop:screen_identification} and~\ref{prop:wb_ml_identification}). Both can be verified by computing the rank of the stacked matrix (see Section~\ref{sec:mc:check_identification}). In practice, we expect both to hold for heterogeneous observed social networks, but emphasise that the relevant condition should be verified for the chosen estimator.

    \item \textbf{Social media use.} For all the individuals of the network, we need to measure their SMU\@. This measure must be quantitative (e.g., daily screen time), and needs to contain some variation to identify $\beta_1$ and to make the CoMO variable $c_i = \max(0, \sbar_i - s_i)$ meaningful. While self-reported measures are possible, they can introduce measurement errors. Preferably, researchers would have available direct device-level usage data to avoid this.

    \item \textbf{Well-being outcomes.} Lastly, we need a measure of WB for the individuals. As with SMU, this should be a quantitative measure with some variation in the sample. Ideally, this measure would also not be based on a single self-reported item, but rather on a comprehensive validated psychometric scale. The collection of this measure should be contemporaneous with both the SMU data and network data to match the static, cross-sectional structure of the model.
\end{enumerate}

Perhaps the most natural or easy-to-collect source for this type of data would be school-based surveys. These could collect both SMU and WB data, while using classrooms and friendship nominations to define the network. An entire school of pupils might almost constitute a full network, as connections to pupils from other schools may be somewhat rare. This should be verified before application. The practical challenge is to gather all three at the same time for the same set of individuals in a way that captures all, or most, of the network structure.

\subsection{Estimation Strategy in Practice}
\label{sec:discussion_conclusion:estimation_strategy}
Along with the data requirements above, we provide two practical recommendations following the results of the simulations in Chapters~\ref{chap:monte_carlo} and~\ref{chap:simultaneous}.

The first recommendation is to treat the ML estimator as primary for the recursive CoMO model when the relevant networks are dense. We believe a network based on the typical classroom qualifies as dense in this context. This recommendation stems from the density sensitivity study of Section~\ref{sec:mc:density}, which showed that the 2SLS estimator exhibits weak-instrument noise under dense networks, while the ML estimator only degrades modestly. The density sensitivity of the FIML estimator has not yet been studied. For dense networks, the safest approach is therefore to use the recursive ML estimator. As a diagnostic test for weak instruments, we recommend reporting the first-stage $F$-statistic. When the first-stage $F$-statistic is below the \textcite[p.~557]{Staiger1997} threshold of $10$, caution is warranted. The Hausman test, which we also ran, should not be relied upon at moderate sample sizes due to its low power under these settings (see Section~\ref{sec:mc:overid}). We additionally recommend fitting both estimators and comparing them. If the estimates disagree and the network is dense with low first-stage $F$ values, the ML estimate should be treated as the central point estimate, and the 2SLS discrepancy as a sign of weak-instrument bias.

Second, when choosing between the recursive CoMO model and the simultaneous CoMO model, we recommend treating the recursive model as the baseline and the simultaneous model as a robustness extension. The baseline recommendation comes from the extensive validation of Chapter~\ref{chap:monte_carlo}, including the density study, and from the analytical standard errors that allow robust inference. If feedback from WB to SMU is expected, or the application is sensitive to the possibility, the simultaneous model should be fitted. The Wald test for $\beta_2 = 0$ can be used to provide evidence about feedback, but only has moderate power at the sample sizes studied in Section~\ref{sec:simul:wald_test}. Thus, non-rejection should not be taken as strong evidence against feedback. When both models are fitted, which is recommended for robustness, agreement between estimates should strengthen confidence in the results. Disagreement should be seen as evidence that the recursive assumption might not hold.

%% Section 8.6: Limitations and Future Research
\section{Limitations and Future Research}
\label{sec:discussion_conclusion:limitations_future}
This section discusses limitations of the thesis, along with the directions of future work that could address them. The main limitation of the thesis is that the modelling framework has not yet been applied to real data. This is clearly the next step for model validation. However, there are also various further limitations to the framework itself. We split these into three categories: assumptions and definitions made within the current analytical framework that could be refined, assumptions made that would require new data or modelling strategies if relaxed, and computational limitations. For each of these we identify the limitations and discuss potential future research avenues.

\subsection{Refinements Within the Current Framework}

The recursive CoMO model makes multiple assumptions that could be relaxed under the same overall modelling strategy.

We begin by mentioning the \emph{recursive assumption} on the structural form, which is already relaxed in Chapter~\ref{chap:simultaneous}. The discussion of Section~\ref{sec:simul:future} mentions limitations to the current work on the simultaneous CoMO model: we have no analytical Hessian, broader Monte Carlo validation with a density-sensitivity study is needed, instrumental variable (IV) methods have not been developed for estimation, and sharper equilibrium conditions could be derived. All of these are items for future work that would validate the estimator, and allow for wider and more robust application of the model. In particular, the development of IV methods would relax the normality assumption, increasing robustness at the cost of some efficiency.

The recursive CoMO model also uses the \emph{normality assumption} (Assumption~\ref{ass:normality}). This is needed for the identification of the peer conformity parameter in the SMU equation via Proposition~\ref{prop:screen_identification} when there are no exogenous covariates present. It also underpins the efficiency of the ML estimator. Our Monte Carlo simulations validated the estimators under a DGP conforming to this assumption, and did not examine robustness to non-normality. Future work could examine how the model performs under heavy-tailed or skewed errors, and develop quasi-ML estimators with sandwich-type inference robust to these scenarios.

The normality assumption also relates to the bounds on SMU and WB in our setup, as the normal errors are unbounded. In practice, this makes the recursive and simultaneous likelihoods local linear approximations in areas away from the boundary. Our Monte Carlo validation is performed at values avoiding this, but a pile-up of values at the boundary (e.g., zero SMU, or ceiling effects in WB) would require model alterations, potentially with a specification including censoring or truncation.

A third assumption made in the recursive CoMO model is the specific form of the CoMO mechanism as a term with the max operator. This creates a sharp kink at $s_i = \sbar_i$, which is arguably not realistic for the potentially more gradual process of social comparison. While other alternatives, such as a smooth soft-hinge $\log(1 + \exp(\sbar_i - s_i))$, were discussed in Section~\ref{sec:model:discussion}, we did not implement or test sensitivity to other choices of the functional form of CoMO\@.

Similarly, the direct effect $\gamma_1 s_i$ of own SMU on WB is imposed as linear, but could in theory have a different functional form. Non-linearities such as saturation or threshold effects cannot be modelled by the current approach. A simple change of the functional form to, for example, a spline in $s_i$, would open the door to more complex relationships. Alternative specifications have not been tested.

Another natural extension within the same framework would be the introduction of heterogeneous peer effects. For example, $\beta_1$ or $\gamma_3$ could be allowed to vary across individuals to add realism.

Finally, the Monte Carlo validation of Chapter~\ref{chap:monte_carlo} is limited in the form of network generation. All networks used for validation are generated from the Erd\H{o}s--R\'enyi model (Section~\ref{sec:mc:network}). While any network satisfying the identification and regularity conditions of Section~\ref{sec:model:identification} will lead to valid estimators, finite-sample behaviour can vary. Model validation with various network architectures, including more realistic topologies, is thus an important direction for future work.

\subsection{Assumptions Requiring New Data or Modelling Strategies}

The models built in this thesis also have more fundamental limitations that cannot be solved by adjusting the current framework alone.

The strongest of these is perhaps the assumption of \emph{network exogeneity} (Assumption~\ref{ass:network_exog}). We believe this is the strongest, as it is almost certainly violated to some degree in real social networks. The literature shows that homophily, the tendency of similar individuals to form connections \parencite{McPherson2001}, is often present in social networks. If connections are made partly due to similarity, there might be unobserved factors driving both SMU and WB, leading to confounding in our model. Addressing this presents a difficulty as future research would need to either explicitly model the network formation process itself, or include covariates strong enough to absorb the sources of homophily. Following the route of additional covariates might also allow 2SLS estimation of the SMU equation, providing a robustness check currently unavailable in the parsimonious specification of Chapter~\ref{chap:model}.

A related limitation is the assumption of \emph{error independence} (Assumption~\ref{ass:error_indep}). This requires $\vect{u} \perp \vect{\varepsilon}$, which would be violated by the presence of, e.g., personality traits or life events that simultaneously affect SMU and WB\@. Further work is needed to relax this assumption.

A further limitation of the models of this thesis is that they are static and cross-sectional, only able to capture a single time point. In Chapter~\ref{chap:abm} we used ABM simulations with temporal dynamics to illustrate Simpson's Paradox. Capturing such effects is outside the scope of the current modelling framework. Modelling dynamic temporal network effects is potentially very policy-relevant, and thus represents an important area of future research. Doing this would require an expansion of the current DGP\@. This might be done through a panel extension allowing lagged SMU and WB, enabling the application of longitudinal data. A starting point could be the spatial-panel machinery of \textcite{LeeYu2010}. This would also allow for modelling of the \emph{cost of having missed out} introduced by \textcite{DahlFoldnesGronneberg2026}, which represents the long-run accumulated WB cost of deprivation of interaction with peers during formative years. Extensions in this direction can also allow for applications of estimators to synthetic data generated from ABMs, along with robustness tests to the misspecification implied by the ABM/analytical gap.

Lastly, the current model setup is limited by requiring a single fully observed network. In reality, individuals are part of multiple networks that overlap to varying degrees (school, family, and online networks). These might also vary in their relevance for social comparison. Future work that deals explicitly with missing edges, false positives, strength of social connections, and other variations in $\mat{G}$ would alleviate this problem.

\subsection{Computational Constraints}

The last category of limitations we will discuss concerns model scalability.

The models developed in the thesis contain terms that are expensive to evaluate for large $n$. In the analytical Hessian for the WB equation, the trace term $\tr[(\Ay^{-1}\mat{G})^2]$ has a computational complexity of $O(n^3)$. In the simultaneous CoMO model, the $\log|\det(\mat{J})|$ term in the likelihood is also expensive, although our Schur-complement simplification in terms of the sparse matrices $\Ay$, $\As$, and $\mat{\Gamma}$ reduces the cost from $O(n^3)$ to one that scales with the non-zero entries of these matrices. In practice, our models are therefore applicable to networks of up to a few thousand individuals. National-scale surveys or other large datasets would be prohibitively large for our models. This also has consequences for model validation, which is expensive to perform for large $n$ and becomes intractable when $n$ grows past a few thousand. Future research could approach this by developing less computationally expensive methods, such as applying the stochastic trace estimator outlined in Section~\ref{sec:est:computational} and using sparse-matrix computations.

%% Section 8.7: Concluding Remarks
\section{Concluding Remarks}
\label{sec:discussion_conclusion:closing}
We conclude by returning to the motivating question of the thesis. The answer to the question is that the causal influence of SMU on WB remains unclear. This thesis contributes not by settling the question, but by clarifying one reason why the existing evidence cannot do so: network effects. We develop the formal machinery for taking network effects into account, allowing for a better answer to be obtained in the future.

This contribution should be seen in light of the limitations discussed in Section~\ref{sec:discussion_conclusion:limitations_future}. Some of the implications, such as the insufficiency of individual-level interventions, are also conditional on empirical validation of the proposed CoMO mechanism.

Despite these caveats, we believe the framework developed in this thesis provides the statistical foundation needed to move past a methodological blind spot in the SMU--WB field: assuming away the very feature that makes social media social.
